\section{Conclusion}
\label{sec:conclusion}

We introduced CIPHER, a procedurally generated benchmark for video-editing captcha solving
that stress-tests VLMs on a compound task: parsing natural-language edit instructions,
executing an ffmpeg transformation chain, and extracting a hidden key from the result.
The three-stage pipeline structure and PipeF1 metric enable fine-grained diagnosis
of where agent capability breaks down.

Our experiments reveal that frontier VLMs solve S3 key extraction near-perfectly
in isolation (Opus/Sonnet: 100\% EM), while OCR baselines struggle ($\sim$25\%).
The compound difficulty of CIPHER lies in the full agentic pipeline:
temporal key localization, correct operation sequencing, and tool execution
create a challenging integration problem that no evaluated model solves reliably
across all levels.

\paragraph{Limitations.}
CIPHER focuses on seven deterministic ffmpeg operations;
generative edits (e.g., inpainting, style transfer) and audio-only keys are out of scope.
Human baselines are not included in this version.

\paragraph{Future work.}
L6 (QR/morse key encoding), multi-agent solving, and adversarial instruction generation
are natural extensions. The procedural generator enables zero-cost dataset refresh,
making CIPHER suitable for ongoing leaderboard tracking.

\paragraph{Reproducibility.}
All code and generation scripts are released at
\url{https://github.com/Fr0do/cipher}.
The test manifest (seeds 0--49) is pinned at commit \texttt{TODO}.
